\documentclass[12pt]{article}
\usepackage[a4paper,margin=1in]{geometry}
\usepackage[utf8]{inputenc}
\usepackage[T1]{fontenc}
\usepackage{amsmath,amssymb}
\usepackage{textcomp}
\usepackage{parskip}
\begin{document}
\begin{center}
{\Large \textbf{1999 STEP III}}
\end{center}

\subsection*{Page 1}
\begin{flushleft}
STEP III, 1999 2\par
\bigskip\textbf{Section A: Pure Mathematics}\par
\bigskip\textbf{Question 1.} Consider the cubic equation\par
3 2\par
\[x px + qx r = 0 ,\]
- -\par
\[where p = 0 and r = 0.\]
6 6\par
\textbf{(i)} If the three roots can be written in the form ak\par
-1,\par
a and ak for some constants a\par
\[and k, show that one root is q/p and that q^3 rp3 = 0 .\]
-\par
\textbf{(ii)} If r = q3/p3 , show that q/p is a root and that the product of the other two roots is\par
(q/p)2. Deduce that the roots are in geometric progression.\par
\textbf{(iii)} Find a necessary and sufficient condition involving p, q and r for the roots to be in\par
arithmetic progression.\par
\bigskip\textbf{Question 2.} (i) Let f (x) = (1 + x2)ex. Show that f 0 (x) > 0 and sketch the graph of f (x). Hence, or\par
otherwise, show that the equation\par
\[(1 + x 2 )ex = k,\]
\[where k is a constant, has exactly one real root if k > 0 and no real roots if k 6 0.\]
\textbf{(ii)} Determine the number of real roots of the equation\par
\[(ex 1) k tan -1 x = 0\]
- -\par
\[in the cases (a) 0 < k 6 2/\pi and (b) 2/\pi < k < 1.\]
3 Justify, by means of a sketch, the formula\par
n 2\par
\[lim f (1 + m/n) = f (x) dx .\]
n\to\infty\par
(\par
n\par
) 1\par
\[m=1 Z\]
X\par
Show that\par
1 1 1\par
\[lim + + + = ln 2 .\]
n\to\infty n + 1 n + 2 \cdots n + n\par
½ ¾\par
Evaluate\par
n n n\par
lim + + + .\par
\[n\to\infty n^2 + 1 n^2 + 4 \cdots n^2 + n^2\]
½ ¾\par
\end{flushleft}
\subsection*{Page 2}
\begin{flushleft}
STEP III, 1999 3\par
\bigskip\textbf{Question 4.} A polyhedron is a solid bounded by F plane faces, which meet in E edges and V vertices. You\par
\[may assume Euler's formula, that V E + F = 2.\]
-\par
In a regular polyhedron the faces are equal regular m-sided polygons, n of which meet at each\par
vertex. Show that\par
4n\par
\[F = ,\]
h\par
\[where h = 4 (n 2)(m 2).\]
- - -\par
By considering the possible values of h, or otherwise, prove that there are only five regular\par
polyhedra, and find V , E and F for each.\par
\bigskip\textbf{Question 5.} The sequence u0, u1, u2, ... is defined by\par
\[u^0 = 1, u^1 = 1, u n+1 = u n + u n-1 for n > 1 .\]
Prove that\par
2 2 2 2\par
\[u n+2 + u n-1 = 2(u n+1 + u n ) .\]
Using induction, or otherwise, prove the following result:\par
2 2 2 2\par
\[u2n = u n + u n-1 and u2n+1 = u n+1 - u n-1\]
for any positive integer n.\par
\bigskip\textbf{Question 6.} A closed curve is given by the equation\par
x\par
2/n\par
+ y\par
2/n\par
\[= a\]
2/n\par
( )\par
∗\par
\[where n is an odd integer and a is a positive constant. Find a parametrization x = x(t),\]
\[y = y(t) which describes the curve anticlockwise as t ranges from 0 to 2\pi.\]
\[Sketch the curve in the case n = 3, justifying the main features of your sketch.\]
The area A enclosed by such a curve is given by the formula\par
\[2\pi\]
dy(t) dx(t)\par
\[A = x(t) y(t) dt .\]
2 dt - dt\par
Z · ¸\par
\[Use this result to find the area enclosed by ( ) for n = 3.\]
∗\par
\end{flushleft}
\subsection*{Page 3}
\begin{flushleft}
STEP III, 1999 4\par
\bigskip\textbf{Question 7.} Let a be a non-zero real number and define a binary operation on the set of real numbers by\par
\[x y = x + y + axy .\]
∗\par
Show that the operation is associative.\par
∗\par
Show that (G, ) is a group, where G is the set of all real numbers except for one number\par
∗\par
which you should identify.\par
Find a subgroup of (G, ) which has exactly 2 elements.\par
∗\par
\bigskip\textbf{Question 8.} The function y(x) is defined for x > 0 and satisfies the conditions\par
dy\par
\[y = 0 and = 1 at x = 0.\]
dx\par
When x is in the range 2(n 1)\pi < x < 2n\pi, where n is a positive integer, y(t) satisfies the\par
-\par
differential equation\par
\[d2y\]
\[+ n y = 0.\]
\[dx2\]
dy\par
\[Both y and are continuous at x = 2n\pi for n = 0, 1, 2, . . . .\]
dx\par
\textbf{(i)} Find y(x) for 0 6 x 6 2\pi.\par
\textbf{(ii)} Show that y(x) = 1 sin 2x for 2\pi 6 x 6 4\pi, and find y(x) for all x > 0.\par
\textbf{(iii)} Show that\par
\infty \infty\par
\[y dx = \pi .\]
\[n^2\]
\[Z n=1\]
X\par
\end{flushleft}
\subsection*{Page 4}
\begin{flushleft}
STEP III, 1999 5\par
\bigskip\textbf{Section B: Mechanics}\par
\bigskip\textbf{Question 9.} The gravitational force between two point particles of masses m and m 0 is mutually attractive\par
and has magnitude\par
Gmm\par
,\par
\[r^2\]
where G is a constant and r is the distance between them.\par
A particle of unit mass lies on the axis of a thin uniform circular ring of radius r and mass m,\par
at a distance x from its centre. Explain why the net force on the particle is directed towards\par
the centre of the ring and show that its magnitude is\par
Gmx\par
.\par
\[(x^2 + r^2)3/2\]
The particle now lies inside a thin hollow spherical shell of uniform density, mass M and radius\par
a, at a distance b from its centre. Show that the particle experiences no gravitational force\par
due to the shell.\par
\bigskip\textbf{Question 10.} A chain of mass m and length l is composed of n small smooth links. It is suspended vertically\par
over a horizontal table with its end just touching the table, and released so that it collapses\par
inelastically onto the table. Calculate the change in momentum of the (k + 1)th link from the\par
bottom of the chain as it falls onto the table.\par
Write down an expression for the total impulse sustained by the table in this way from the\par
whole chain. By approximating the sum by an integral, show that this total impulse is\par
approximately\par
\[m\sqrt(2gl)\]
when n is large.\par
11 Calculate the moment of inertia of a uniform thin circular hoop of mass m and radius a about\par
an axis perpendicular to the plane of the hoop through a point on its circumference.\par
The hoop, which is rough, rolls with speed v on a rough horizontal table straight towards the\par
edge and rolls over the edge without initially losing contact with the edge. Show that the\par
hoop will lose contact with the edge when it has rotated about the edge of the table through\par
an angle θ, where\par
\[1 v^2\]
\[cos θ = + .\]
2 2ag\par
\end{flushleft}
\subsection*{Page 5}
\begin{flushleft}
STEP III, 1999 6\par
\bigskip\textbf{Section C: Probability and Statistics}\par
\bigskip\textbf{Question 12.} In the game of endless cricket the scores X and Y of the two sides are such that\par
(j + k)λj+k\par
-1\par
\[P(X = j, Y = k) = e ,\]
j!k!\par
\[for some positive constant λ, where j, k = 0, 1, 2, . . . .\]
\textbf{(i)} Find P(X + Y = n) for each n > 0.\par
\textbf{(ii)} Show that\par
\[2λe2λ-1\]
\[= 1.\]
\textbf{(iii)} Show that\par
\[2xe2x-1\]
\[is an increasing function of x for x > 0 and deduce that the\]
equation in (ii) has at most one solution and hence determine λ.\par
\textbf{(iv)} Calculate the expectation E(2X+Y ).\par
\bigskip\textbf{Question 13.} The cakes in our canteen each contain exactly four currants, each currant being randomly\par
placed in the cake. I take a proportion X of a cake where X is a random variable with density\par
function\par
\[f(x) = Ax\]
for 0 6 x 6 1 where A is a constant.\par
\textbf{(i)} What is the expected number of currants in my portion?\par
\textbf{(ii)} If I find all four currants in my portion, what is the probability that I took more than\par
half the cake?\par
\bigskip\textbf{Question 14.} In the basic version of Horizons (H1) the player has a maximum of n turns, where n > 1. At\par
each turn, she has a probability p of success, where 0 < p < 1. If her first success is at the rth\par
turn, where 1 6 r 6 n, she collects r pounds and then withdraws from the game. Otherwise,\par
\[her winnings are nil. Show that in H^1, her expected winnings are\]
p -1 1 + nqn+1 (n + 1)qn pounds,\par
-\par
\[where q = 1 p. \pounds  ¤\]
-\par
The rules of H2 are the same as those of H1, except that n is randomly selected from a Poisson\par
\[distribution with parameter λ. If n = 0 her winnings are nil. Otherwise she plays H^1 with\]
\[the selected n. Show that in H^2, her expected winnings are\]
1 e\par
-λp\par
λqe\par
-λp\par
pounds.\par
p - -\par
³ ´\par
\end{flushleft}
\end{document}
